\documentclass[11pt]{article}

\usepackage{amsmath}
\usepackage{amsfonts}
\usepackage{graphicx}
\usepackage{hyperref}
\usepackage{physics}

\title{Computer Vision Algorithms Used in the MedVision AI Project}
\author{Yann Sofian Smatti}
\date{}

\begin{document}

\maketitle

\tableofcontents

\newpage

\section{Introduction}

Computer Vision aims to allow machines to understand visual data.
In the medical domain, this includes analyzing radiography, MRI, CT scans,
and other diagnostic images.

In this project we rely mainly on:

\begin{itemize}
\item Convolutional Neural Networks
\item Transfer Learning
\item Image Preprocessing
\item Explainable AI
\end{itemize}

\section{Mathematical Representation of Images}

A digital image can be represented as a matrix.

For a grayscale image:

\[
I \in \mathbb{R}^{H \times W}
\]

where

\begin{itemize}
\item $H$ = height
\item $W$ = width
\end{itemize}

Each element corresponds to the pixel intensity.

\[
I(x,y) \in [0,255]
\]

For RGB images:

\[
I \in \mathbb{R}^{H \times W \times 3}
\]

\section{Image Normalization}

Neural networks perform better when data is normalized.

Normalization transforms the pixel values:

\[
x' = \frac{x}{255}
\]

so that

\[
x' \in [0,1]
\]

This stabilizes gradient descent.

\section{Convolution Operation}

Convolution is the core operation in CNNs.

Given an image $I$ and a kernel $K$:

\[
(I * K)(x,y) = \sum_{i=-k}^{k} \sum_{j=-k}^{k} I(x-i,y-j)K(i,j)
\]

The kernel extracts features such as edges.

Example edge detection kernel:

\[
K =
\begin{bmatrix}
1 & 0 & -1 \\
1 & 0 & -1 \\
1 & 0 & -1
\end{bmatrix}
\]

This operation produces a \textbf{feature map}.

\section{Convolutional Neural Networks}

CNNs consist of stacked layers.

Typical architecture:

\begin{itemize}
\item Convolution
\item Activation
\item Pooling
\item Fully connected layers
\end{itemize}

A convolutional layer computes:

\[
h = f(W * x + b)
\]

where

\begin{itemize}
\item $W$ = convolution filter
\item $b$ = bias
\item $f$ = activation function
\end{itemize}

\section{Activation Functions}

The most common activation function is ReLU:

\[
ReLU(x) = \max(0,x)
\]

Advantages:

\begin{itemize}
\item avoids vanishing gradients
\item faster training
\end{itemize}

\section{Pooling Layers}

Pooling reduces spatial resolution.

Max pooling:

\[
y = \max(x_1,x_2,x_3,x_4)
\]

Example:

\[
\begin{bmatrix}
1 & 3 \\
2 & 0
\end{bmatrix}
\rightarrow 3
\]

Benefits:

\begin{itemize}
\item reduces computation
\item increases translation invariance
\end{itemize}

\section{Fully Connected Layers}

After feature extraction, features are flattened:

\[
x \in \mathbb{R}^{n}
\]

Prediction is obtained via:

\[
y = \sigma(Wx + b)
\]

where $\sigma$ is the sigmoid function.

\section{Binary Cross Entropy Loss}

For binary classification:

\[
L = - \frac{1}{N} \sum_{i=1}^{N}
\left[y_i \log(p_i) + (1-y_i)\log(1-p_i)\right]
\]

where

\begin{itemize}
\item $y_i$ = ground truth
\item $p_i$ = predicted probability
\end{itemize}

\section{Gradient Descent}

Parameters are updated using gradient descent.

\[
\theta_{t+1} = \theta_t - \eta \nabla L(\theta)
\]

where

\begin{itemize}
\item $\eta$ = learning rate
\item $\nabla L$ = gradient of the loss
\end{itemize}

\section{Backpropagation}

Backpropagation computes gradients using the chain rule.

For a composition of functions:

\[
f(x) = f_3(f_2(f_1(x)))
\]

the gradient becomes:

\[
\frac{df}{dx} =
\frac{df_3}{df_2}
\cdot
\frac{df_2}{df_1}
\cdot
\frac{df_1}{dx}
\]

This allows efficient training of deep networks.

\section{Transfer Learning}

Training deep CNNs from scratch requires millions of images.

Transfer learning uses pretrained models such as:

\begin{itemize}
\item ResNet
\item EfficientNet
\item VGG
\end{itemize}

Mathematically:

\[
f(x) = g(h(x))
\]

where

\begin{itemize}
\item $h(x)$ = pretrained feature extractor
\item $g(x)$ = new classifier
\end{itemize}

Only the last layers are trained.

\section{Data Augmentation}

Medical datasets are often small.

Data augmentation artificially increases the dataset.

Transformations include:

\begin{itemize}
\item rotation
\item flipping
\item zoom
\item brightness change
\end{itemize}

Mathematically:

\[
x' = T(x)
\]

where $T$ is a transformation.

\section{Explainable AI}

In medical applications, interpretability is critical.

Grad-CAM generates heatmaps highlighting important regions.

\[
L^c = ReLU\left(
\sum_k \alpha_k^c A^k
\right)
\]

where

\begin{itemize}
\item $A^k$ = feature maps
\item $\alpha_k^c$ = weights computed via gradients
\end{itemize}

The result is a localization map.

\section{Evaluation Metrics}

Accuracy:

\[
Accuracy =
\frac{TP + TN}{TP + TN + FP + FN}
\]

Precision:

\[
Precision =
\frac{TP}{TP + FP}
\]

Recall (Sensitivity):

\[
Recall =
\frac{TP}{TP + FN}
\]

F1 score:

\[
F1 =
2\frac{Precision \cdot Recall}{Precision + Recall}
\]

\section{ROC Curve}

The ROC curve represents the tradeoff between
True Positive Rate and False Positive Rate.

\[
TPR =
\frac{TP}{TP + FN}
\]

\[
FPR =
\frac{FP}{FP + TN}
\]

The Area Under Curve (AUC) measures classifier performance.

\section{Conclusion}

Computer Vision techniques combined with deep learning allow automated medical image analysis.

However, for clinical use, models must be:

\begin{itemize}
\item robust
\item explainable
\item validated on large datasets
\end{itemize}

\end{document}